\documentclass[11pt]{amsart}

\usepackage[a4paper,margin=1in]{geometry}
\usepackage[T1]{fontenc}
\usepackage{lmodern}
\usepackage{microtype}
\usepackage{amsmath,amssymb,amsthm,mathtools}

\numberwithin{equation}{section}

\newtheorem{theorem}{Theorem}[section]
\newtheorem{lemma}[theorem]{Lemma}
\newtheorem{definition}[theorem]{Definition}
\newtheorem{remark}[theorem]{Remark}

\title{Chebyshevization of the Real-Axis Singular-Value Problem}
\author{}

\begin{document}

\maketitle

\section{Chebyshevization of the fourth-order recurrence}

Fix
\[
        0<a<1,
        \qquad
        A_n(a)=\operatorname{tridiag}(a,0,1),
\]
and put
\[
        B_n(x):=xI_n-A_n(a),
        \qquad
        P_n(x):=B_n(x)^TB_n(x).
\]
Since
\[
        s_{\min}(B_n(x))^2
        =
        \lambda_{\min}(P_n(x)),
\]
the real-axis singular-value problem for \(B_n(x)\) is equivalent to the
ordinary eigenvalue problem
\[
        P_n(x)v=\eta v,
        \qquad \eta=s^2.
\]

\subsection{The pentadiagonal matrix \(P_n(x)\)}

The matrix \(B_n(x)\) has diagonal entries \(x\), superdiagonal entries
\(-1\), and subdiagonal entries \(-a\).  Therefore \(P_n(x)=B_n(x)^TB_n(x)\)
is real symmetric and pentadiagonal.  Its entries are
\[
        (P_n(x))_{j,j}
        =
        \begin{cases}
        x^2+a^2, & j=1,\\
        x^2+1+a^2, & 2\le j\le n-1,\\
        x^2+1, & j=n,
        \end{cases}
\]
\[
        (P_n(x))_{j,j+1}
        =
        (P_n(x))_{j+1,j}
        =
        -(1+a)x,
        \qquad 1\le j\le n-1,
\]
and
\[
        (P_n(x))_{j,j+2}
        =
        (P_n(x))_{j+2,j}
        =
        a,
        \qquad 1\le j\le n-2.
\]

Thus, in the interior, the eigenvalue equation
\[
        P_n(x)v=\eta v
\]
has the constant-coefficient fourth-order form
\begin{equation}\label{eq:fourth-order}
        a v_{j+2}
        -(1+a)xv_{j+1}
        +(x^2+1+a^2-\eta)v_j
        -(1+a)xv_{j-1}
        +a v_{j-2}
        =
        0 .
\end{equation}

The endpoint corrections can be encoded by ghost boundary conditions.  The
left endpoint equations are equivalent to
\begin{equation}\label{eq:left-ghost}
        v_0=0,
        \qquad
        av_{-1}+v_1=0,
\end{equation}
and the right endpoint equations are equivalent to
\begin{equation}\label{eq:right-ghost}
        v_{n+1}=0,
        \qquad
        v_{n+2}+av_n=0.
\end{equation}
With these ghost conditions, the same fourth-order recurrence
\eqref{eq:fourth-order} holds formally for \(j=1,\dots,n\).

\begin{lemma}[Ghost boundary formulation]
The equation
\[
        (P_n(x)-\eta I_n)v=0
\]
is equivalent to the fourth-order recurrence
\[
        a v_{j+2}
        -(1+a)xv_{j+1}
        +(x^2+1+a^2-\eta)v_j
        -(1+a)xv_{j-1}
        +a v_{j-2}
        =
        0,
        \qquad 1\le j\le n,
\]
together with the boundary conditions
\[
        v_0=0,\qquad av_{-1}+v_1=0,
\]
and
\[
        v_{n+1}=0,\qquad v_{n+2}+av_n=0.
\]
\end{lemma}

\begin{proof}
For \(2\le j\le n-1\), the \(j\)-th row of
\[
        (P_n(x)-\eta I_n)v=0
\]
is exactly \eqref{eq:fourth-order}.  At \(j=1\), the actual first row is
\[
        (x^2+a^2-\eta)v_1-(1+a)xv_2+av_3=0.
\]
The formal fourth-order equation at \(j=1\) is
\[
        av_3-(1+a)xv_2+(x^2+1+a^2-\eta)v_1-(1+a)xv_0+av_{-1}=0.
\]
Using
\[
        v_0=0,
        \qquad
        av_{-1}+v_1=0,
\]
this reduces exactly to the true first-row equation.

Similarly, at \(j=n\), the actual last row is
\[
        av_{n-2}-(1+a)xv_{n-1}+(x^2+1-\eta)v_n=0.
\]
The formal equation at \(j=n\) contains the additional terms
\[
        a v_{n+2}-(1+a)xv_{n+1}+a^2v_n.
\]
Using
\[
        v_{n+1}=0,
        \qquad
        v_{n+2}+av_n=0,
\]
these extra terms vanish.  Hence the ghost formulation is equivalent to the
original finite matrix problem.
\end{proof}

\subsection{The palindromic characteristic polynomial}

Substituting the trial solution
\[
        v_j=z^j
\]
into the fourth-order recurrence \eqref{eq:fourth-order} gives
\[
        a z^{j+2}
        -(1+a)xz^{j+1}
        +(x^2+1+a^2-\eta)z^j
        -(1+a)xz^{j-1}
        +az^{j-2}
        =
        0.
\]
After division by \(z^{j-2}\), one obtains the characteristic polynomial
\begin{equation}\label{eq:palindromic-quartic}
        a z^4
        -(1+a)xz^3
        +(x^2+1+a^2-\eta)z^2
        -(1+a)xz
        +a
        =
        0 .
\end{equation}

This polynomial is palindromic: the coefficient of \(z^4\) equals the
constant coefficient, and the coefficient of \(z^3\) equals the coefficient
of \(z\).  Equivalently, after dividing by \(z^2\), it is invariant under
\[
        z\longmapsto z^{-1}.
\]
Indeed, division by \(z^2\) gives
\[
        a(z^2+z^{-2})
        -(1+a)x(z+z^{-1})
        +(x^2+1+a^2-\eta)
        =
        0.
\]

Now introduce the palindromic variable
\[
        u:=z+z^{-1}.
\]
Since
\[
        z^2+z^{-2}
        =
        (z+z^{-1})^2-2
        =
        u^2-2,
\]
the quartic equation \eqref{eq:palindromic-quartic} reduces to the quadratic
equation
\begin{equation}\label{eq:u-quadratic}
        a u^2-(1+a)x u+x^2+(1-a)^2-\eta=0.
\end{equation}

Thus the fourth-order problem splits into two second-order channels.  Let
\(u_+(x,\eta)\) and \(u_-(x,\eta)\) be the two roots of
\eqref{eq:u-quadratic}:
\begin{equation}\label{eq:u-pm}
        u_\pm(x,\eta)
        =
        \frac{
        (1+a)x
        \pm
        \sqrt{(1-a)^2(x^2-4a)+4a\eta}
        }{2a}.
\end{equation}
Then the fourth-order operator factors as
\begin{equation}\label{eq:operator-factorization}
        a\bigl(E+E^{-1}-u_+\bigr)
         \bigl(E+E^{-1}-u_-\bigr),
\end{equation}
where \(E\) is the shift operator
\[
        (Ef)_j=f_{j+1}.
\]

\begin{lemma}[Palindromic factorization]
The fourth-order recurrence \eqref{eq:fourth-order} factors as
\[
        a\bigl(E+E^{-1}-u_+\bigr)
         \bigl(E+E^{-1}-u_-\bigr)v=0,
\]
where \(u_\pm\) are given by \eqref{eq:u-pm}.  Equivalently, every solution
is built from the two second-order Chebyshev channels
\[
        f_{j+1}=u_+ f_j-f_{j-1},
        \qquad
        f_{j+1}=u_- f_j-f_{j-1}.
\]
\end{lemma}

\begin{proof}
The sum and product of the two roots of \eqref{eq:u-quadratic} are
\[
        u_++u_-=\frac{(1+a)x}{a},
\]
and
\[
        u_+u_-=\frac{x^2+(1-a)^2-\eta}{a}.
\]
Therefore
\[
\begin{aligned}
&a(E+E^{-1}-u_+)(E+E^{-1}-u_-)                         \\
&=
a(E+E^{-1})^2
-a(u_++u_-)(E+E^{-1})
+au_+u_-                                                 \\
&=
a(E^2+2+E^{-2})
-(1+a)x(E+E^{-1})
+x^2+(1-a)^2-\eta                                        \\
&=
aE^2-(1+a)xE+(x^2+1+a^2-\eta)
-(1+a)xE^{-1}+aE^{-2}.
\end{aligned}
\]
This is precisely the recurrence operator in \eqref{eq:fourth-order}.
\end{proof}

\subsection{Chebyshev channels}

For \(j\ge0\), define
\[
        C_j(u):=T_j\!\left(\frac u2\right),
\]
and
\[
        S_j(u):=U_{j-1}\!\left(\frac u2\right),
        \qquad S_0(u):=0,\quad S_1(u):=1.
\]
Here \(T_j\) and \(U_j\) are the Chebyshev polynomials of the first and second
kind.  Thus
\[
        C_0(u)=1,
        \qquad
        C_1(u)=\frac u2,
\]
and both sequences satisfy the same second-order recurrence
\[
        F_{j+1}(u)=uF_j(u)-F_{j-1}(u).
\]
In particular,
\[
        C_{j+1}(u)=uC_j(u)-C_{j-1}(u),
\]
and
\[
        S_{j+1}(u)=uS_j(u)-S_{j-1}(u).
\]

The sequence \(C_j(u)\) is even in the index:
\[
        C_{-j}(u)=C_j(u),
\]
while the sequence \(S_j(u)\), extended by the same recurrence to negative
indices, satisfies
\[
        S_{-j}(u)=-S_j(u).
\]
In particular,
\[
        S_{-1}(u)=-1.
\]

When \(u=z+z^{-1}\), these functions have the elementary representations
\[
        C_j(u)=\frac{z^j+z^{-j}}2,
\]
and
\[
        S_j(u)=\frac{z^j-z^{-j}}{z-z^{-1}},
\]
with the usual limiting interpretation when \(z=\pm1\).

\begin{lemma}[Chebyshev basis for each channel]
For each root \(u_\sigma\), \(\sigma\in\{+,-\}\), the two sequences
\[
        C_j(u_\sigma),
        \qquad
        S_j(u_\sigma)
\]
solve the second-order channel equation
\[
        f_{j+1}=u_\sigma f_j-f_{j-1}.
\]
Consequently, they also solve the fourth-order recurrence
\eqref{eq:fourth-order}.
\end{lemma}

\begin{proof}
The recurrence relation is one of the defining recurrence relations for
Chebyshev polynomials:
\[
        T_{j+1}(y)=2yT_j(y)-T_{j-1}(y),
\]
and
\[
        U_j(y)=2yU_{j-1}(y)-U_{j-2}(y).
\]
With \(y=u/2\), both become
\[
        F_{j+1}=uF_j-F_{j-1}.
\]
Since the fourth-order operator factors as
\[
        a(E+E^{-1}-u_+)(E+E^{-1}-u_-),
\]
any solution of either second-order channel is automatically a solution of
the fourth-order equation.
\end{proof}

When
\[
        u_+\ne u_-,
\]
the four sequences
\[
        C_j(u_+),\quad S_j(u_+),\quad C_j(u_-),\quad S_j(u_-)
\]
span the four-dimensional solution space of the fourth-order recurrence.
Thus the general bulk solution is
\begin{equation}\label{eq:general-cheb-solution}
        v_j
        =
        \alpha_+ C_j(u_+)
        +\beta_+ S_j(u_+)
        +\alpha_- C_j(u_-)
        +\beta_- S_j(u_-).
\end{equation}

If \(u_+=u_-\), the four displayed sequences are no longer linearly
independent.  In that case the correct basis is obtained by taking the
divided-difference limit, equivalently by adjoining
\[
        \partial_u C_j(u),
        \qquad
        \partial_u S_j(u).
\]
All determinant formulas below are interpreted by analytic continuation
through the repeated-channel set \(u_+=u_-\).

\subsection{Boundary determinant in Chebyshev form}

The finite eigenvalue problem is obtained by imposing the four ghost boundary
conditions
\[
        v_0=0,
        \qquad
        av_{-1}+v_1=0,
\]
and
\[
        v_{n+1}=0,
        \qquad
        v_{n+2}+av_n=0.
\]
Using the Chebyshev representation \eqref{eq:general-cheb-solution}, these
four boundary conditions give a homogeneous linear system in the four
coefficients
\[
        \alpha_+,\ \beta_+,\ \alpha_-,\ \beta_-.
\]

Define the \(4\times4\) Chebyshev boundary matrix.  In the display below,
write \(C_k^\pm=C_k(u_\pm)\) and \(S_k^\pm=S_k(u_\pm)\):
\[
        \mathcal B_n(x,\eta)
        :=
        \begin{bmatrix}
        C_0^+ & S_0^+ & C_0^- & S_0^-\\[0.3em]
        aC_{-1}^+ + C_1^+
        &
        aS_{-1}^+ + S_1^+
        &
        aC_{-1}^- + C_1^-
        &
        aS_{-1}^- + S_1^-\\[0.3em]
        C_{n+1}^+ & S_{n+1}^+ & C_{n+1}^- & S_{n+1}^-\\[0.3em]
        C_{n+2}^+ + aC_n^+
        &
        S_{n+2}^+ + aS_n^+
        &
        C_{n+2}^- + aC_n^-
        &
        S_{n+2}^- + aS_n^-
        \end{bmatrix}.
\]
Since
\[
        C_0(u)=1,\qquad S_0(u)=0,
\]
and
\[
        C_{-1}(u)=C_1(u)=\frac u2,
        \qquad
        S_{-1}(u)=-1,\qquad S_1(u)=1,
\]
the first two rows may also be written as
\[
        [\,1,\ 0,\ 1,\ 0\,],
\]
and
\[
        \left[
        \frac{1+a}{2}u_+,\ 1-a,\ \frac{1+a}{2}u_-,\ 1-a
        \right].
\]

Define the Chebyshev boundary determinant
\[
        \mathfrak D_n(x,\eta):=\det \mathcal B_n(x,\eta).
\]

\begin{theorem}[Chebyshev boundary determinant]
Assume first that \(u_+\ne u_-\).  Then
\[
        \eta\in\sigma(P_n(x))
        \quad\Longleftrightarrow\quad
        \mathfrak D_n(x,\eta)=0.
\]
Equivalently,
\[
        \mathfrak D_n(x,\eta)=0
\]
is the Chebyshev form of the finite eigenvalue condition
\[
        \det(P_n(x)-\eta I_n)=0.
\]
At points where \(u_+=u_-\), the same equivalence holds after interpreting
\(\mathfrak D_n\) by divided-difference continuation.
\end{theorem}

\begin{proof}
When \(u_+\ne u_-\), every solution of the fourth-order recurrence has the
form
\[
        v_j
        =
        \alpha_+ C_j(u_+)
        +\beta_+ S_j(u_+)
        +\alpha_- C_j(u_-)
        +\beta_- S_j(u_-).
\]
Substituting this expression into the four ghost boundary conditions gives
\[
        \mathcal B_n(x,\eta)
        \begin{bmatrix}
        \alpha_+\\ \beta_+\\ \alpha_-\\ \beta_-
        \end{bmatrix}
        =
        0.
\]
A nonzero vector \(v\) satisfying the finite eigenvalue equation exists if
and only if this homogeneous system has a nontrivial solution, which is
equivalent to
\[
        \det\mathcal B_n(x,\eta)=0.
\]
The repeated-channel case follows by analytic continuation, or equivalently
by replacing the repeated columns by the corresponding derivatives with
respect to \(u\).
\end{proof}

\subsection{Meaning of Chebyshevization}

The preceding construction is what we call the \emph{Chebyshevization} of the
real-axis singular-value problem.  It consists of the following steps:

\begin{enumerate}
\item Start from the singular-value problem
\[
        s_{\min}(xI_n-A_n(a))^2
        =
        \lambda_{\min}\bigl((xI_n-A_n(a))^T(xI_n-A_n(a))\bigr).
\]

\item Write the eigenvalue equation
\[
        P_n(x)v=\eta v
\]
as the scalar fourth-order constant-coefficient recurrence
\[
        a v_{j+2}
        -(1+a)xv_{j+1}
        +(x^2+1+a^2-\eta)v_j
        -(1+a)xv_{j-1}
        +a v_{j-2}=0.
\]

\item Observe that the characteristic polynomial
\[
        a z^4
        -(1+a)xz^3
        +(x^2+1+a^2-\eta)z^2
        -(1+a)xz
        +a
\]
is palindromic.

\item Use the palindromic substitution
\[
        u=z+z^{-1}
\]
to reduce the quartic characteristic equation to the quadratic equation
\[
        a u^2-(1+a)xu+x^2+(1-a)^2-\eta=0.
\]

\item Let \(u_+\) and \(u_-\) be the two roots of this quadratic.  The
fourth-order recurrence then factors into two second-order Chebyshev
channels:
\[
        f_{j+1}=u_+ f_j-f_{j-1},
        \qquad
        f_{j+1}=u_- f_j-f_{j-1}.
\]

\item Express all bulk solutions by Chebyshev polynomials:
\[
        T_j\!\left(\frac{u_\pm}{2}\right),
        \qquad
        U_j\!\left(\frac{u_\pm}{2}\right).
\]

\item Finally, impose the finite ghost boundary conditions to obtain an
explicit Chebyshev boundary determinant
\[
        \mathfrak D_n(x,\eta)=0.
\]
\end{enumerate}

Thus the original real-axis singular-value problem for the nonnormal
tridiagonal Toeplitz matrix has been converted into a two-channel Chebyshev
boundary problem.

\end{document}
